\section{Limitations and Threats to Validity}\label{sec:limits}

The findings of Sections~\ref{sec:background} to~\ref{sec:discussion} are only as strong as the measurements behind
them. This section says where they are weak, starting with the weakness that matters most.

\paragraph{One boot per experiment.}
The controlled evidence comes from three pod sessions on three different hosts (\dv{lean.session.minutes},
\dv{lean.session2.minutes} and \dv{lean.session3.minutes} minutes, Table~\ref{tab:pod-hardware}). Each experiment's
measurements come from a single boot. The \texttt{gen30} configuration was booted on all three pods, but no
measurement was repeated across boots, so the intervals cover repetitions inside a boot and nothing across boots. The
one configuration measured on two hosts differed by far more than any interval inside a boot
(Section~\ref{sec:results-controlled}), which makes the placement of a pod a large factor that we did not control. E8
has two repetitions per level in scenario B, but in scenario A only at $N=1$, because one client-side disconnect among
eight requests tripped the pre-registered stop rule (error rate above 5\%) and ended A's second sweep. A rule evaluated
on eight requests is too sensitive. After the first session we closed idle client connections before the server does,
which was the cause of those disconnects, and required at least two errors, a rule that E12 and E13 used. The second
session's rerun of the E3 subset kept the first session's rule to match its baseline, and the first session's data
were not rerun under the new rule.

\paragraph{The questions the campaign leaves open.}
Our three sessions (Section~\ref{sec:setup}) answer RQ1 on the measured hardware, give RQ2 two host shares with one
repetition each, measure the RQ4 frontier for two scenarios, give RQ5 a first answer on small subsets, and answer RQ6
for one pod and six context regimes with one repetition (E13, Section~\ref{sec:agent}). They leave open the following.
\begin{itemize}
\item The speed of the model held entirely in GPU memory, so that without the 2x B300 pod (E9) the cost of the offload
is bounded only by the slope between two host shares, not against a baseline without offload.
\item RQ3 as a comparison of speed, because without Diffbot's hooks the zero-copy variant cannot serve DeepSeek-V4.1 on
these pins at any KV block size (64 fails at initialization, 128 on the first request), so the cost of the per-expert
split against upstream's offload is unknown.
\item A quality reference on our own hardware and serving code, since the API stands in for E11 and differences may
come from the API's precision, prompt handling or sampling rather than from quantization.
\item Upstream vLLM with the original checkpoint on the same GPUs, dropped to avoid its \dv{static.orig.total_gib}\,GiB download.
\item Plans fitted to routing counts at several shares and the smallest share that fits, and the decode scenario at the
larger share, lost to a client-side disconnect.
\item Placement policy (E4), Engram format (E5), the SM budget of the host-expert stream (E6) and speculation off (E7).
\item The prose scenario, Poisson arrivals, MATH, MBPP, RULER and LongBench, while long context (E12) was measured with
one repetition, on the slower second pod, and without a quality check at those lengths.
\end{itemize}
Statements about these remain hypotheses.

\paragraph{Model capability.}
The LiveBench scores in Section~\ref{sec:background} belong to the unquantized model at maximum reasoning effort as
served by DeepSeek. How much of that capability survives \dv{cfg.exl3_bits}-bit quantization with thinking disabled is
not measured, since our quality evidence covers three small task subsets with thinking off. Running LiveBench on the
served stack, with thinking on and off, is future work.

\paragraph{The coding-agent case study.}
E13 is one repetition on one boot of one pod. Levels near the objective pass or fail by noise, which is why we
provision on a 10\% margin, and its capacities do not transfer to another pod (Section~\ref{sec:results-controlled}).
Replies have a fixed length of 400 tokens, and the agent turns are synthetic, with unique corpus text standing in for
tool output, so speculative decoding may accept more or fewer tokens than on real coding work. The usage profile is
one developer's week, read from a provider's billing categories that need not match vLLM's prefix-cache accounting, and
the reply length per turn and the daily active window are assumptions varied in the tables. The duty cycle uses the
floor of 40 tokens/s, an upper bound on busy time, and the Binomial model assumes independent developers with equal
activity, which understates the tail for a team on one schedule. Costs are GPU rental at one quoted price and exchange
rate. Storage, networking, engineering, operations and the differences in model capability and product features are
not priced, except in the stated sensitivity.

\paragraph{The exploratory data.}
The imported runs were measured once per configuration, used usable bars chosen after seeing the data, started every
level with a synchronized burst that dominates the TTFT tails, depend on content through speculative decoding with
synthetic prompts (slices of the vLLM sources and question templates, or a few essays), and never evaluated quality.
They also timed chunks rather than tokens, computed percentiles by nearest rank on 8 to 128 requests, measured
aggregate throughput including ramp-up and drain, and ran on one pod on one day.

\paragraph{The static accounting.}
File bytes are not runtime allocations. Repacking, FP8 copies and workspaces change the memory used on the GPU
(\dv{expl.runtime_minus_file_gib}\,GiB per GPU more in the exploratory boot). The placement rules model the documented
configuration and are checked against the server log only where the log reports placement. The minimal host share is
an estimate from one boot.

\paragraph{The protocol.}
The service-level objectives were registered before the controlled runs but after we saw the exploratory data. Cloud
pods share hosts, so CPU, memory bandwidth and PCIe topology vary between pods. We record them but cannot control
them, and repetitions across boots were designed for one pod and did not run. Upstream vLLM's zero-copy offload did
fail on this model at the pinned commit, at both KV block sizes we tried, and we report that failure as a result, but
it leaves RQ3 unanswered as a comparison of speed. When the B300 reference runs, it will bound the cost of the offload
together with the difference in hardware, because it uses other GPUs. GSM8K, HumanEval and similar benchmarks may be
in the training data, and since the comparison is between two versions of the same model, that limits the concern but
does not remove it. NVML's PCIe counters are sampled and their meaning varies between GPUs, hence the calibration
step.

\paragraph{Attribution.}
The audit relies on pinned public sources, keyword searches confirmed by reading code, and metadata from the Hub and
GitHub. Diffbot's code is attributed to the organization, because individual authors are not visible. Work that is
not public or not cited by the recipe could precede it.
